\documentclass[10pt]{notes}
 
\begin{document}

\thispagestyle{firstpage}
\begin{center}
    \large\textsc{Simulation Ideas}
\end{center}

For each observation $i\in [n]$, denote by its tensor covariate $\mathcal{X}_i\in\bbR^{p_1\times\cdots\times p_D}$, one-way covariates $\mathbf{z}_i\in\bbR^{p_0}$, and two responses $y_{ij}\in\bbR$ where $j=1,2$ represents two eyes. Let $t\in[0,1]$ be the observation time, $\mathcal{A}(t)=(\alpha_{j_1,\dots,j_d}(t))\in\bbR^{p_1\times\cdots\times p_D}$ be the unknown tensor coefficient function, and $\bb\in\bbR^{p_0}$ be the coefficients. The model is assumed to be
\begin{align*}
y_{ij}
&= \left\langle \mathcal{X}_{ij}, \mathcal{A}(t_i) \right\rangle
   + \mathbf{z}_i^\top \boldsymbol{\beta}
   + \epsilon_{ij} \\
&= \sum_{s=1}^S
   \left\langle \mathcal{X}_{ij}^{(s)}, \mathcal{A}^{(s)}(t_i) \right\rangle
   + \mathbf{z}_i^\top \boldsymbol{\beta}
   + \epsilon_{ij} \\
\intertext{by a rank-$R$ CP decomposition,}
&\approx \sum_{s=1}^S \sum_{r=1}^R
   x_{ijrs}^{*}
   \left\langle
   \mathcal{U}_r^{(s)},
   \mathcal{A}^{(s)}(t_i)
   \right\rangle
   + \mathbf{z}_i^\top \boldsymbol{\beta}
   + \epsilon_{ij} \\
&= \left\langle \mathbf{X}_{ij}^{*}, \mathbf{A}(t_i) \right\rangle
   + \mathbf{z}_i^\top \boldsymbol{\beta}
   + \epsilon_{ij},
\end{align*}
where matrix $\mathbf{X}_{ij}^{*}\coloneq(x_{ijrs}^*)\in\bbR^{R\times S}$, and $\mathbf{A}(t)\coloneq (\alpha_{rs}(t))\in\bbR^{R\times S}$ with $\alpha_{rs}(t)\coloneq \innerprod{\mathcal{U}_r^{(s)}}{\mathcal{A}^{(s)}(t)}$. Within one patient, the eyes are correlated 
\begin{align*}
    \begin{bmatrix}
        \epsilon_{i1}\\ \epsilon_{i2}
    \end{bmatrix}\sim 
    \mvnorm\begin{pmatrix}
        \mathbf{0},\begin{bmatrix}
            \sigma^2 & \rho\sigma^2\\ 
            \rho\sigma^2 & \sigma^2
        \end{bmatrix}
    \end{pmatrix}.
\end{align*}
\textcolor{red}{\textbf{Remark.} In the current Cases I--IV below, we use a simplified single-eye setting and do not consider within-subject correlation between the two eyes: no correlation in $(\epsilon_{i1},\epsilon_{i2})^\top$. Hence, the DGP needs to be updated, and the method codes also needs to be updated accordingly.}

\subsubsection*{Current Case I \& II}
The tensor size is set to be $D=3$, $p_1=p_2=p_3=40$, and $S=64$. We split the tensor by $S^{(1)}=S^{(2)}=S^{(3)}=10$ and $p_1'=p_2'=p_3'=4$. Set the rank $R=5$, then the whole-sample $s$-th sub-block tensor $\tilde{\mathcal{X}}^{(s)}\in\bbR^{n\times p_1'\times p_2'\times p_3'}$ can be generated by 
\begin{align*}
    \tilde{\mathcal{X}}^{(s)}=\sum_{r=1}^{R} \mathbf{x}_{rs}^*\circ \mathbf{x}_r^{(s1)}\circ \mathbf{x}_r^{(s2)}\circ \mathbf{x}_r^{(s3)}
\end{align*} 
where $\mathbf{x}_{rs}^*\in\bbR^n$ is a Gaussian ensemble vector whose entries are iid from $\normal(0,1)$, and any $\mathbf{x}_{r}^{(sd)}$ and $\mathbf{x}_{r'}^{(sd)}$ are orthonormal and obtained from the QR decomposition of a random matrix. The coefficient tensor function is set to be
\begin{align*}
    \mathcal{A}^{(s)}(t)=\sum_{r=1}^{R} \alpha_{rs}(t)\mathbf{x}_r^{(s1)}\circ \mathbf{x}_r^{(s2)}\circ \mathbf{x}_r^{(s3)}
\end{align*}
where $\alpha_{rs}(t)$ has 4 choises: $4\sqrt{\frac{rs}{RS}}(t-0.5)^2, \sqrt{\frac{rs}{RS}}\, t^{0.5}, 1.75\sqrt{\frac{rs}{RS}}
\left(\exp\{-(3t-1)^2\} + \exp\{-(4t-3)^2\} - 0.75\right)$,
and
$\sqrt{\frac{rs}{RS}} \sin\bigl(2\pi(t-0.5)\bigr)$.
The one-way vector $\mathbf{Z}\in\bbR^{n\times p_0}$ is Gaussian ensemble matrix, and $\bb=(3,3)^\top$.

\begin{remark}
    Case I is just the 2D version of Case II.
\end{remark}

\subsubsection*{Current Case III \& IV}

For Case IV, the tensor size is set to be $D=3$, $p_1=p_2=p_3=80$, and
$S=64$. We split the tensor by
$S^{(1)}=S^{(2)}=S^{(3)}=4$, so that
$p_1'=p_2'=p_3'=20$. Set the rank $R=20$, then the whole-sample $s$-th
sub-block tensor
$\tilde{\mathcal{X}}^{(s)} \in \mathbb{R}^{n \times p_1' \times p_2' \times p_3'}$
can be generated by
$
    \tilde{\mathcal{X}}^{(s)}
    =
    \sum_{r=1}^{R}
    \mathbf{x}_{rs}^*
    \circ
    \mathbf{x}_r^{(s1)}
    \circ
    \mathbf{x}_r^{(s2)}
    \circ
    \mathbf{x}_r^{(s3)},
$
where $\mathbf{x}_{rs}^* \in \mathbb{R}^n$ is a Gaussian ensemble vector
whose entries are iid from $\normal(0,1)$, and any
$\mathbf{x}_{r}^{(sd)}$ and $\mathbf{x}_{r'}^{(sd)}$ are orthonormal and
obtained from the QR decomposition of a random matrix. The coefficient
tensor function is set to be
$
    \mathcal{A}^{(s)}(t)
    =
    \sum_{r=1}^{R}
    \alpha_{rs}(t)
    \mathbf{x}_r^{(s1)}
    \circ
    \mathbf{x}_r^{(s2)}
    \circ
    \mathbf{x}_r^{(s3)},
$
where
\begin{align*}
    \alpha_{rs}(t)
    =
    \begin{cases}
        \sqrt{\frac{rs}{RS}} \sin\bigl(2\pi(t-0.5)\bigr),
        & 1 \le r \le s \le 10, \\[4pt]
        \sqrt{\frac{rs}{RS}},
        & 1 \le s < r \le 10, \\[4pt]
        0,
        & \text{otherwise}.
    \end{cases}
\end{align*}
The one-way vector $\mathbf{Z} \in \mathbb{R}^{n \times p_0}$ is a Gaussian
ensemble matrix, and $\bb = (1,1,0,0,0)^\top$. The random error
$\epsilon_i$ follows $\normal(0,1)$. This case is designed for the
high-dimensional setting, where $2RS + p_0 > n$, so the initial local linear
kernel smoothing step is no longer feasible and the paper focuses on the
penalized spline stage and the refined regression stage.

\begin{remark}
    Case III is the 2D counterpart of Case IV, with spatial correlation across partitions being
    introduced. More specifically, for two partitions indexed by
    $s=(s^{(1)},s^{(2)})$ and $s'=(s'^{(1)},s'^{(2)})$, the corresponding
    latent scores are correlated according to an AR(1)-type structure with
    covariance
    \[
    \rho^{\,|s^{(1)}-s'^{(1)}| + |s^{(2)}-s'^{(2)}|}.
    \]
    Hence Case III is used to assess structure identification under spatially
    correlated tensor covariates, whereas Case IV is used to assess the same
    problem under a high-dimensional 3D setting.
\end{remark}

\subsection*{Next Step 1: Two-Eye Update}
This is the most urgent update. We can keep the signal structure in Cases I--IV unchanged, while extending the data-generating process to incorporate a two-eye structure: correlation in $(\epsilon_{i1},\epsilon_{i2})^\top$. The corresponding estimation procedure also needs to be updated so that subject-level dependence between the two eyes is properly accounted for.

\subsection{Other Idea: Generating Tensor Covariate}
Currently, the tensor $\tilde{\mathcal{X}}$ is generated by tensor product of vector, like 
$\tilde{\mathcal{X}}^{(s)}
    =
    \sum_{r=1}^{R}
    \mathbf{x}_{rs}^*
    \circ
    \mathbf{x}_r^{(s1)}
    \circ
    \mathbf{x}_r^{(s2)}
    \circ
    \mathbf{x}_r^{(s3)}$.
To make the revised simulation more convincing, we can directly start generating covariate tensor $\mathcal{X}_i$ and the coefficient tensor $\mathcal{A}(t)$, so that the CP decomposition really has some information loss.

\subsubsection{By Matrix Normal}
In the 2D case, we may generate
\[
\mathrm{vec}(\mathcal{X}_i)
\sim
N\!\left(
\mathbf{0},
\Sigma_2 \otimes \Sigma_1
\right),
\]
where $\mathcal{X}_i \in \mathbb{R}^{p_1 \times p_2}$ and each
$\Sigma_d \in \mathbb{R}^{p_d \times p_d}$. This is the matrix-normal
construction, and its higher-order extension leads to the tensor-normal model
with separable covariance structure; see \cite{manceur2013}. A simple choice is
$
(\Sigma_d)_{uv} = \exp\!\left(- |u-v|/\ell_d \right),
$
or a Mat\'ern-type covariance if we want more flexible local regularity. In
the 3D case, we may analogously generate
\[
\mathrm{vec}(\mathcal{X}_i)
\sim
N\!\left(
\mathbf{0},
\Sigma_3 \otimes \Sigma_2 \otimes \Sigma_1
\right),
\]
for $\mathcal{X}_i \in \mathbb{R}^{p_1 \times p_2 \times p_3}$. This gives a
directly generated image-like tensor with spatial correlation.

\subsubsection{By Gaussian Process}
We can also
generally generate $\mathcal{X}_i$ as a spatial Gaussian random field or
Gaussian process. In the 2D case, for pixel locations
$s=(u,v)$ and $s'=(u',v')$, we may specify
\[
\mathcal{X}_i(s) \sim \mathcal{GP}\bigl(0, C(s,s')\bigr),
\]
where the covariance kernel can be chosen as
\[
C(s,s')
=
\tau^2 \exp\left\{
-\frac{(u-u')^2}{\ell_1^2}
-\frac{(v-v')^2}{\ell_2^2}
\right\},
\]
or more generally as a Mat\'ern kernel. 3D case can mimic this idea.
For a general review of Gaussian random field generation methods, see
\cite{liu2019}. For an image-regression example where predictor images are
simulated from Gaussian processes, see \cite{guo2022}.

\subsubsection{By Real Images Empirical Basis}

Another possibility is to use real images to construct a more realistic
simulation mechanism. Suppose we have preprocessed training images from the
real-data study. Let $\bar{\mathcal{X}}$ be the sample mean image and let
$\Psi_1,\dots,\Psi_{K_0}$ be the first several empirical principal component
images extracted from the observed tensors. Then we may generate
\[
\mathcal{X}_i
=
\bar{\mathcal{X}}
+ \sum_{k=1}^{K_0} \xi_{ik}\Psi_k
+ \mathcal{E}_i,
\]
where the scores $\xi_{ik}$ are sampled from a fitted distribution, for
example $\xi_{ik}\sim \normal(0,\lambda_k)$ with $\lambda_k$ equal to the
estimated eigenvalue of the $k$-th principal component, and
$\mathcal{E}_i$ is a residual spatial noise field. This empirical-PCA or
eigenimage construction is close in spirit to simulation designs built from
real neuroimaging scans; see \cite{goldsmith2014}.

\subsection{Generating Tensor Coefficient Function}
The coefficient tensor can be generated by
$
\mathcal{A}(t)
=
\sum_{k=1}^K g_k(t)\,\Phi_k,
$
where each $\Phi_k$ is a spatial template on the original image domain and
each $g_k(t)$ controls how the strength of that region varies with $t$. For the 3D case, we may define
three fixed spatial templates by
\begin{align*}
    \Phi_1(u,v,w)
    &= \exp\!\left\{
    -\frac{(u-d_{11})^2 + (v-d_{12})^2 + (w-d_{13})^2}{2h_1^2}
    \right\}, \\
    \Phi_2(u,v,w)
    &= \exp\!\left\{
    -\left(
    \frac{(u-d_{21})^2}{2a_1^2}
    + \frac{(v-d_{22})^2}{2a_2^2}
    + \frac{(w-d_{23})^2}{2a_3^2}
    \right)
    \right\}, \\
    \Phi_3(u,v,w)
    &= -\exp\!\left\{
    -\frac{(u-d_{31})^2 + (v-d_{32})^2 + (w-d_{33})^2}{2h_3^2}
    \right\},
\end{align*}
where $\Phi_1$ and $\Phi_3$ are two localized spherical templates centered at
$(d_{11},d_{12},d_{13})$ and $(d_{31},d_{32},d_{33})$, and $\Phi_2$ is an
elongated ellipsoidal region centered at $(d_{21},d_{22},d_{23})$. 2D case can be generated in the similar way.

But I am still considering how to add sparsity in this setting. Maybe we can still adopt the current $\mathcal{A}(t)$ generation method in Case IV.

\begin{thebibliography}{9}

\bibitem{manceur2013}
Manceur, A. M. and Dutilleul, P. (2013).
Maximum likelihood estimation for the tensor normal distribution:
Algorithm, minimum sample size, and empirical bias and dispersion.
\textit{Journal of Computational and Applied Mathematics}, \textbf{239}, 37--49.

\bibitem{liu2019}
Liu, Y., Li, J., Sun, S., and Yu, B. (2019).
Advances in Gaussian random field generation: a review.
\textit{Computational Geosciences}, \textbf{23}(5), 1011--1047.

\bibitem{guo2022}
Guo, C., Kang, J., and Johnson, T. D. (2022).
A spatial Bayesian latent factor model for image-on-image regression.
\textit{Biometrics}, \textbf{78}(1), 72--84.

\bibitem{goldsmith2014}
Goldsmith, J., Huang, L., and Crainiceanu, C. M. (2014).
Smooth scalar-on-image regression via spatial Bayesian variable selection.
\textit{Journal of Computational and Graphical Statistics}, \textbf{23}(1), 46--64.

\end{thebibliography}


\end{document}
